\documentclass[10pt,a4paper]{amsart}
\usepackage[margin=19mm]{geometry}
\usepackage{amsmath,amssymb,mathtools}
\usepackage[scr=rsfs,cal=euler]{mathalfa}
\usepackage{xfrac}
\usepackage[unicode,hidelinks,bookmarks=false]{hyperref}
\newcommand{\ie}{i.e.\ }
\newcommand{\cf}{cf.\ }
\newcommand{\resp}{respectively\ }
\newcommand{\Subsection}{\subsection}
\providecommand{\nobreakdash}{\nobreak\hbox{-}}
\setlength{\emergencystretch}{2em}
\multlinegap0pt
\usepackage{bm}
\usepackage{bbm}
\usepackage{dsfont}
\usepackage{xspace}
\usepackage{cancel}
\usepackage{mathtools}
\usepackage[shortlabels]{enumitem}
\usepackage{amsthm}
\makeatletter
\@ifundefined{corollary}{\newtheorem{corollary}{Corollary}}{}
\@ifundefined{definition}{\newtheorem{definition}{Definition}}{}
\@ifundefined{theorem}{\newtheorem{theorem}{Theorem}}{}
\makeatother
\DeclarePairedDelimiter{\norm}{\lVert}{\rVert}
\newcommand{\N}{\mathbb{N}}
\newcommand{\Oc}{\mathcal{O}}
\newcommand{\R}{\mathbb{R}}
\newcommand{\p}{\partial}
\title[High-order integration on regular triangulated manifolds rea]{High-order integration on regular triangulated manifolds reaches super-algebraic approximation rates through cubical re-parameterizations}
\author{Gentian Zavalani, Oliver Sander, Michael Hecht}
\date{Source: arXiv:2311.13909v4 (CC BY 4.0)}
\begin{document}
\maketitle

\noindent\emph{Source and licence.} High-order integration on regular triangulated manifolds reaches super-algebraic approximation rates through cubical re-parameterizations. Version arXiv:2311.13909v4 (CC BY 4.0), distributed under the Creative Commons Attribution 4.0 International licence, as recorded in the arXiv licence metadata for this exact version and shown on the abstract page https://arxiv.org/abs/2311.13909v4. Full rights notice: licenses/arxiv-2311.13909-CC-BY-4.0.txt.\par\smallskip

\noindent\textit{Editorial scope.} Counted unit: the Corollary whose source label is \texttt{coir:trunc} in the article (source0.tex lines 1027--1045, source label \texttt{coir:trunc}), delivered together with the article's own complete original proof of it (source0.tex lines 1046--1125, one \texttt{proof} environment of 80 lines). No mathematics is written, repaired or re-derived: statement and proof are the article's own text line for line. Required context retained uncounted: def:total\_variation (lines 462--470); theo:TRUNC (lines 986--998); eq:bound (lines 993--996). Every definition, display and label of those lines is retained unchanged. Omitted from this package and not redistributed: 1--461, 471--985, 997--1026, 1126--2417. Nothing inside a retained excerpt is omitted or altered: every retained body is delivered exactly as the article's lines are written, so no omission receipt of any kind accompanies this package. Citations: the retained excerpts carry no citation command at all, so no bibliography entry of the article is transcribed and no reference is redistributed. Reference adaptation: none was needed and none was made; every reference inside the delivered unit resolves inside it, so no unresolved reference or citation is produced. Printed slips kept literal: no printed word, symbol, hypothesis or number of the counted statement or of its proof was altered. Layout and class: the article's own class is \texttt{siamart220329} with packages including listings, lipsum, amsfonts, graphicx, epstopdf, algorithmic, natbib, amsopn, amsmath, amssymb, mathtools, mathrsfs, mathptmx, booktabs; the wrapper is the lane's pinned \texttt{amsart} editorial wrapper at 10pt on A4, which restates the theorem-like environments the retained text needs and adds the article's own macro definitions that the retained text needs (\texttt{\textbackslash N}, \texttt{\textbackslash Oc}, \texttt{\textbackslash R}, \texttt{\textbackslash norm}, \texttt{\textbackslash p}), with extra wrapper packages (bm, bbm, dsfont, xspace, cancel, mathtools, enumitem). The delivered numbering is the unit's own. Assets: no retained span contains a figure environment, a graphics-inclusion command, a PDF-page inclusion, a table or a TikZ picture; no image file is copied into this package, the package declares no asset dependency and no image file is redistributed with it. Licence: version arXiv:2311.13909v4 of this article is distributed under the Creative Commons Attribution 4.0 International licence, as recorded in the arXiv licence metadata for this exact version and shown on the abstract page https://arxiv.org/abs/2311.13909v4. A full rights notice is delivered at licenses/arxiv-2311.13909-CC-BY-4.0.txt. Characters: every retained excerpt is pure ASCII, so the delivered engine, XeLaTeX with its default Latin Modern text font, reports no missing character.
\begin{definition}[$r^\text{th}$-order total variation]
\label{def:total_variation}
Let $r \geq 0$, $f: \square_d \to \R$ and its derivatives through $\p^{\boldsymbol{\beta}} f$, $\boldsymbol{\beta} = (r+1,\dots,r+1)$ be absolutely continuous (differentiable almost everywhere). We define the $r^{\text{th}}$ total variation $V_{f,r}$ as
\begin{equation}\label{eq:var}
   V_{f,r}=\max_{\substack{\boldsymbol{\beta} \in \N^d \\ \|\boldsymbol{\beta}\|_\infty \leq r+1}}\int_{\square_d} |\p^{\boldsymbol{\beta}}f(\mathrm{x})| d\mathrm{x} \,, 
   % \quad \omega(\mathrm{x}) = \prod_{i=1}^d \frac{1}{\sqrt{1-x_i^2}} \,,
\end{equation}
and refer $f$ as having bounded $r^{\text{th}}$ total variation, whenever $V_{f,r} < \infty$ exists.
\end{definition}

 \begin{theorem}\label{theo:TRUNC}  Let $d\in\N$, $r\geq 0$, and $f$ be of bounded $r^{\text{th}}$ total variation, Definition~\ref{def:total_variation}. 
 % multivariate function with absolutely continuous derivatives up to order
% $r-1$, possessing bounded $r^{\text{th}}$ total variation $0\leq V_{f,r} <\infty$.
Then $f$ can be expanded in a Chebyshev series
\begin{equation*}
f(\mathrm{x})   =\sum_{\boldsymbol{\alpha}\in \N_0^d}c_{\boldsymbol{\alpha}} T_{\alpha_1}(x_1)\cdots T_{\alpha_d}(x_d)\,,
\end{equation*}
\begin{equation}\label{eq:bound}
    \text{with} \quad
|c_{\boldsymbol{\alpha}}|\leq V_{f,r}\Big(\frac{2}{ \pi q(q-1)\ldots(q-r)}\Big)^{d} \,,
\end{equation}
whenever $q =\min_{i=1,\dots,d} \alpha_i  \geq r+1$.
\end{theorem}

\begin{corollary}\label{coir:trunc} Let the assumptions of Theorem~\ref{theo:TRUNC} be fulfilled. We denote with
\begin{equation}
\mathcal{T}_{f,n}(\mathrm{x})=\sum_{\boldsymbol{\alpha} \in A_{d,n}}c_{\boldsymbol{\alpha}}T_{\alpha_1}(x_1)\cdots T_{\alpha_d}(x_d)
\end{equation}
the truncated Chebyshev series of $f : \square_d \to \R$ with respect to $A_{d,n}$, with $n >r$.
\begin{enumerate}[label=\roman*)]
    \item \label{item:approx_1} The truncation error is bounded by
    \begin{equation}\label{eq:trunc}
        \|f - \mathcal{T}_{f,n}\|_{C^0(\square_d)} \leq \frac{2ed^2V_{f,r}}{\pi(r-d+1)}\left(\frac{n+1}{\,n+1-r\,}\right)^{r+1}\cdot \frac{1}{n^{r-d+1}} \in\Oc\big(n^{-(r-d+1)}\big)\,,
    \end{equation}
    $r>d-1$.
    \item \label{item:approx_2} The truncation 
    error of the first-order partial derivatives is bounded by
        \begin{equation}\label{eq:dtrunc}
        \|\p_{x_i} f - \p_{x_i} \mathcal{T}_{f,n}\|_{C^0(\square_d)} \leq \frac{2ed^2V_{f,r}}{\pi(r-d-1)}\left(\frac{n+1}{\,n+1-r\,}\right)^{r+1}\cdot\frac{1}{n^{r-d-1}}\in\Oc\big(n^{-(r-d-1)}\big)\,, 
    \end{equation}
    $r>d$, $\forall i=1,\dots,d.$
\end{enumerate}
\end{corollary}

\begin{proof}
\ref{item:approx_1} directly follows from Theorem~\ref{theo:TRUNC}:
Since \( T_k(\cos(x)) = \cos(kx) \) for all \( k \in \mathbb{N}_0 \), we observe that \( \|T_k\|_{C^0([-1,1])} \leq 1 \).
Let \(\Lambda_n^{(i)} = \{ \boldsymbol{\alpha} \in \mathbb{N}_{0}^d : \alpha_i > n \}\).
Then the truncation error admits the following bound:

% \ref{item:approx_1} directly follows from Theorem~\ref{theo:TRUNC}:
% Since \( T_k(\cos(x)) = \cos(kx) \) for all \( k \in \mathbb{N} \), we observe that \( \|T_k\|_{C^0([-1,1])} \leq 1 \). Define the set \( \Lambda_n^{(i)} = \left\{ \alpha \in \mathbb{N}^d : \alpha_i > n,\ \alpha_j \leq n \text{ for } j \neq i \right\} \). Then, the truncation error can be estimated as


% \begin{align}
%   \|f - \mathcal{T}_{f,n}\|_{C^0(\Omega_{d})}
%   &\leq
%   \sum_{\boldsymbol{\alpha} \in \mathbb{N}^d \setminus A_{d,n}} c_{\boldsymbol{\alpha}}\| T_{\alpha_1}(x_1)\cdots T_{\alpha_d}(x_d)\|_{C^0(\Omega_{d})}  \leq \sum_{\boldsymbol{\alpha} \in \mathbb{N}^d \setminus A_{d,n}} |c_{\boldsymbol{\alpha}}| \nonumber\\
%   &\textcolor{red}{\leq \sum_{\substack{\alpha \in \mathbb{N}^d \setminus A_{d,n} }}
% \frac{4V_{f,r}}{ \pi(\|\alpha\|_{\infty}-r)^{r+1}}\leq \frac{4V_{f,r}}{\pi}\sum_{k=n+1}^{\infty} \frac{ |A_{d,k}\setminus A_{d,k-1}|}{(k - r)^{r+1}}}\label{eq:conv}\\
%   &\textcolor{red}{= \frac{4V_{f,r}}{\pi}\sum_{k=n+1}^{\infty} \frac{ k^{d-1}}{(k - r)^{r+1}} \leq \frac{4V_{f,r}}{\pi}\int_n^\infty x^{d-r-2} \, dx\nonumber}\\
%   &\textcolor{red}{= \frac{4V_{f,r}}{\pi(d-r-1)}\cdot \frac{1}{n^{(r-d+1)}},}
% \end{align}
% \GZ{Also: For \( d = 1 \), we obtain \(\mathcal{O}(n^{-r})\), which is consistent with Trefethen's theorem.
\begin{align}
  \|f - \mathcal{T}_{f,n}\|_{C^0(\square_{d})}
  &\leq
  \sum_{\boldsymbol{\alpha} \in \mathbb{N}_{0}^d \setminus A_{d,n}} c_{\boldsymbol{\alpha}}\| T_{\alpha_1}(x_1)\cdots T_{\alpha_d}(x_d)\|_{C^0(\Omega_{d})}  \leq \sum_{\boldsymbol{\alpha} \in \mathbb{N}_{0}^d \setminus A_{d,n}} |c_{\boldsymbol{\alpha}}|\nonumber\\
  &\leq \sum_{i=1}^{d}  \sum_{\boldsymbol{\alpha} \in \Lambda_n^{(i)}}
\frac{2V_{f,r}}{ \pi \alpha_i(\alpha_i-1)\ldots(\alpha_i-r)}\leq \sum_{i=1}^{d}  \sum_{\boldsymbol{\alpha} \in \Lambda_n^{(i)}}
\frac{2V_{f,r}}{ \pi (\alpha_i-r)^{r+1}}\nonumber\\
  &= \frac{2dV_{f,r}}{\pi}\sum_{q=n+1}^{\infty} \frac{ |A_{d,q}\setminus A_{d,q-1}|}{(q - r)^{r+1}}\leq \frac{2ed^2V_{f,r}}{\pi}\sum_{q=n+1}^{\infty} \frac{ q^{d-1}}{(q - r)^{r+1}}\nonumber\\
  & \leq \frac{2ed^2V_{f,r}}{\pi}\left(\frac{n+1}{\,n+1-r\,}\right)^{r+1} \sum_{q=n+1}^{\infty} \frac{1}{q^{r-d+2}}\leq\frac{2ed^2V_{f,r}}{\pi}\left(\frac{n+1}{\,n+1-r\,}\right)^{r+1}\int_{n}^{\infty}\frac{1}{x^{r-d+2}}dx\nonumber\\\nonumber
  &= \frac{2ed^2V_{f,r}}{\pi(r-d+1)}\left(\frac{n+1}{\,n+1-r\,}\right)^{r+1}\cdot \frac{1}{n^{r-d+1}}\label{eq:sum_bound}\,
\end{align}
where for each \( \boldsymbol{\alpha} \in \Lambda_n^{(i)} \), we apply the Chebyshev coefficient bound~\eqref{eq:bound} in the \( i \)-th coordinate only, assuming \( q:=\alpha_i \geq r + 1 \). Moreover,we used the estimate \( |A_{d,q} \setminus A_{d,q-1}|
= (q+1)^d - q^d
\;\le\; e d q^{d-1}, \; (q > d)\), and bounded the resulting monotonically decreasing sum (for \(r-d+2 >1\)), which ultimately yields the desired estimate~\eqref{eq:trunc}.







We show \ref{item:approx_2} for the partial derivative $\p_{x_i}$ by writing
\begin{equation*}
    \norm{\partial_{x_i}f-\partial_{x_i}\mathcal{T}_{f,n}}_{C^0(\square_d)}\leq \sum_{\boldsymbol{\alpha} \in \N_0^d \setminus A_{d,n}} |c_{\boldsymbol{\alpha}}|\norm{T_{\alpha_1}\cdots T_{\alpha_{i-1}}}_{C^0(\square_d)}\norm{T'_{\alpha_i}}_{C^0(\square_d)}\norm{T_{\alpha_{i+1}}\cdots T_{\alpha_d}}_{C^0(\square_d)}\,.
\end{equation*}
We recall that $T_{k}(x)=\cos{\left(k\arccos(x\right))}$ for $-1\leq x\leq 1$ and deduce
that for all $k \in \N_0$
\begin{equation}\label{eq.derivative}
T'_{k}(x)=\frac{k\sin{\left(k\arccos\left(x\right)\right)}}{\sqrt{1-x^{2}}} =\frac{k\sin{\left(kt\right)}}{\sin{\left(t\right)}}\,, \qquad t=\arccos(x)\,,
\end{equation}
yielding $\|T'_{\alpha_i}\|_{C^0(\square_d)}=\alpha_i^{2}$.
Following~\ref{item:approx_1}, we compute

\begin{align}
    \norm{\partial_{x_i}f-\partial_{x_i}\mathcal{T}_{f,n}}_{C^0(\square_d)}
    &\leq \sum_{\boldsymbol{\alpha} \in \N_0^d \setminus A_{d,n}} |c_{\boldsymbol{\alpha}}|\alpha_i^2  \leq \sum_{i=1}^{d}  \sum_{\boldsymbol{\alpha} \in \Lambda_n^{(i)}}
 \frac{2V_{f,r}}{ \pi (\alpha_i-r)^{r+1}}\alpha_i^2\nonumber \\ 
 % & \leq \sum_{\alpha \in \mathbb{N}^d \setminus A_{d,n}}
 % V_{f,r}\Big(\frac{2}{ \pi (q-r)^{r+1}}\Big)^{d}\|\alpha\|_\infty^2 \\
&\leq \frac{2ed^2V_{f,r}}{\pi}\sum_{q=n+1}^{\infty} \frac{q^{d-1}q^2}{(q-r)^{r+1}} \leq \frac{2ed^2V_{f,r}}{\pi}\left(\frac{n+1}{\,n+1-r\,}\right)^{r+1} \sum_{q=n+1}^{\infty} \frac{1}{q^{r-d}}\nonumber\\
&\leq \frac{2ed^2V_{f,r}}{\pi}\left(\frac{n+1}{\,n+1-r\,}\right)^{r+1}\int_{n}^{\infty}\frac{1}{x^{r-d}}=\frac{2ed^2V_{f,r}}{\pi(r-d-1)}\left(\frac{n+1}{\,n+1-r\,}\right)^{r+1}\cdot\frac{1}{n^{r-d-1}} \,,\label{eq:EST}
\end{align}
by bounding the monotonically decreasing sum (for  \(r-d >1\))  in \eqref{eq:EST}.
% \red{
% \begin{align}
%     \norm{\partial_{x_i}f-\partial_{x_i}\mathcal{T}_{f,n}}_{C^0(\square_d)}
%     &\leq \sum_{\alpha \in \N^d \setminus A_{d,n}} |c_{\alpha}|\alpha_i^2 \\
%     &\leq \sum_{\substack{\alpha \in \mathbb{N}^d \setminus A_{d,n} \\ q = \min_{i=1,\dots,d} \alpha_i}}
%  V_{f,r}\Big(\frac{2}{ \pi q(q-1)\ldots(q-r)}\Big)^{d} \alpha_i^2\nonumber\\
%     &\leq \sum_{\substack{\alpha \in \mathbb{N}^d \setminus A_{d,n} \\ q = \min_{i=1,\dots,d} \alpha_i}}
%  V_{f,r}\Big(\frac{2}{ \pi q(q-1)\ldots(q-r)}\Big)^{d} \|\alpha\|_\infty^2 \leq \sum_{\substack{\alpha \in \mathbb{N}^d \setminus A_{d,n} \\ q = \min_{i=1,\dots,d} \alpha_i}}
%  V_{f,r}\Big(\frac{2}{ \pi q(q-1)\ldots(q-r)}\Big)^{d} \|\alpha\|_1^2 \nonumber\\
%     & \textcolor{red}{= \frac{2^{d}V_{f,r}}{\pi} \sum_{k=n+1}^{\infty} \frac{k^{d-1}k^2}{k^d(k-1)^d(k-2)^d\cdots(k-r)^d}= \frac{2^{d}V_{f,r}}{\pi} \sum_{k=n+1}^{\infty} \frac{k}{(k-1)^d(k-2)^d\cdots(k-r)^d}}\nonumber\\
%     & \leq \frac{(n+1)2d V_{f,r}}{n^d\pi}\sum_{k=n+1}^{\infty} \frac{1}{(k-2)^d(k-3)^d\cdots(k-r)^d} \label{eq:EST}\\
%     &=\frac{(n+1)2d V_{f,r}}{n^d\pi} \frac{1}{(r-2)^d(n-2)^d(n-3)^d\cdots(n-r)^d}\\
%     &\leq \frac{(n+1)2d V_{f,r}}{n^d\pi} \frac{1}{(r-2)^d(n-r)^{d(r-1)}}\nonumber
% \\    &\textcolor{blue}{\leq \frac{(n+1)2d V_{f,r}}{n^d\pi} \frac{1}{(r-2)^d(n-r)^{(r-1)}}}\nonumber
% \end{align}
% where we have used $\sum_{j=N+1}^{\infty} \frac{1}{j^d(j + 1)^d \cdots (j + m)^d} = \frac{1}{m^d(N + 1)^d(N + 2)^d \cdots (N + m)^d}$ in \eqref{eq:EST}.}
\end{proof}

\end{document}
